\section{Datos, preparación e hipótesis}

La distribución socioeconómica se basa en los subestratos del INSE y la cantidad de viviendas y hogares se aproxima a partir del Censo 2023 \cite{inse2023,censo2023}.

Para Montevideo, la Tabla D3 del INSE 2023 informa las proporciones de hogares B$-$ 5,0\%, B+ 10,4\%, M$-$ 16,0\%, M 21,8\%, M+ 21,2\%, A$-$ 17,1\% y A+ 8,6\%. Estas cifras sustituyen en el escenario base a la distribución preliminar no específica de Montevideo. Los ingresos de referencia continúan siendo los promedios per cápita por subestrato con valor locativo estimados sobre la ECH 2022.

\subsection{Fuentes y transformaciones}

El catálogo de microdatos del INE identifica el estudio como \texttt{URY-INE-CPHV-2023}, con fecha de referencia 31 de mayo de 2023 y unidad geográfica mínima de segmento censal \cite{censo2023ponderado}. La edición de mayo de 2026 incorporó ponderadores para corregir una omisión censal estimada en 10,3\%; en julio de 2026 el INE publicó una revisión posterior del mismo estudio. El catálogo enumera los archivos \texttt{personas\_ext\_05\_2026\_extraccion} y \texttt{viviendas\_ext\_05\_2026\_extraccion}.

El prototipo no descarga, redistribuye ni procesa esos microdatos individuales. Por tanto, no existen archivos censales dentro del repositorio, checksums que registrar ni una licencia que permita presentarlos como parte del proyecto. Se conservan únicamente la identificación y URL del catálogo oficial, cuya consulta está sujeta a los términos del INE: uso científico o estadístico, publicación agregada y prohibición de reidentificación y redistribución sin consentimiento. La implementación usa agregados publicados y parámetros explícitos en archivos de configuración para generar instancias sintéticas reproducibles.

La información municipal de la Intendencia basada en ECH 2023 se utilizará para indicadores socioeconómicos. Sus tablas de cantidades de viviendas corresponden a los censos 2004 y 2011 y no se tratarán como datos de vivienda de 2023.

\subsection{Generación de instancias}

Además del escenario de Montevideo se utilizan instancias sintéticas pequeñas y medianas para verificar reglas e invariantes. Toda generación registra la semilla y los parámetros efectivos.

El generador recorre dos veces la grilla. En la primera pasada determina de forma reproducible cuáles celdas son residenciales y cuáles están ocupadas, obteniendo así la cantidad exacta de hogares para reservar arreglos contiguos. En la segunda construye celdas y hogares, asigna zonas mediante franjas verticales y selecciona subestratos con los pesos configurados. El escenario base usa proporción residencial 0,885 y ocupación 0,90; los conteos efectivos dependen de la semilla y se informan al iniciar cada ejecución.

\subsection{Supuestos y limitaciones de los datos}

No se imputan registros censales ni se realiza una agregación por municipio. Las siete clases socioeconómicas se sortean con las proporciones agregadas para Montevideo del INSE y las zonas del modelo se construyen como franjas verticales de la grilla; estas zonas son particiones computacionales y no municipios reales. Los valores ausentes en las fuentes no se completan: las tolerancias, precios y reglas económicas sin estimación oficial se declaran como supuestos del escenario. Incorporar microdatos ponderados, geometría censal y una correspondencia reproducible entre segmentos y municipios queda fuera del alcance de esta implementación.

La geometría real no se utiliza en el escenario base. Esta convención elimina una dependencia cartográfica, pero impide interpretar posiciones o distancias simuladas como posiciones o distancias reales. La cartografía censal publicada en mayo de 2026, que se encontraba bajo revisión técnica del INE al diseñar el proyecto, queda reservada para una extensión futura.
